\documentclass[11pt]{article}

\usepackage[T1]{fontenc}
\usepackage[utf8]{inputenc}
\usepackage[english]{babel}
\usepackage[top=2.6cm, bottom=2.5cm, left=2.4cm, right=2.4cm]{geometry}
\usepackage{amsmath, amssymb, amsfonts, mathtools, bm, nicefrac}
\usepackage{booktabs}
\usepackage{microtype}
\usepackage{xcolor}
\usepackage{algorithm}
\usepackage{algpseudocode}
\usepackage{fancyhdr}
\usepackage[colorlinks=true, linkcolor=blue, citecolor=blue, urlcolor=blue]{hyperref}

% Matches commands.sty in Tensor_hyper_contraction__THC_, so the bodies are interchangeable.
\newcommand{\mat}[1]{{\mathbf{#1}}}
\newcommand{\parth}[1]{\left(#1\right)}
\newcommand{\sqrbrckts}[1]{\left[#1\right]}

\newcommand{\Dcal}{\mathcal{A}}
\newcommand{\Qcal}{\mathcal{Q}}
\newcommand{\Ncal}{\mathcal{N}}
\newcommand{\eb}{\tilde{\epsilon}}
\newcommand{\Had}{\circ}

\setlength\parindent{0pt}
\setlength\parskip{8pt}

\pagestyle{fancy}
\setlength{\headheight}{14pt}
\lhead{\normalsize\textrm{Marcus Allen}}
\rhead{THC dRPA moments}
\renewcommand{\headrulewidth}{0.4pt}

\title{\vspace{-1.5cm}A separable route to the dRPA density--density moments\\
under tensor hypercontraction}
\author{Marcus Allen}
\date{25 August 2026}

\begin{document}
\maketitle
\vspace{-1.6cm}

\begin{abstract}
\noindent
The direct random-phase approximation supplies the density--density moments that
moment-conserving $GW$ convolves into a self-energy. Under density fitting those moments
are reached by contracting a three-index vertex against the particle--hole energy
denominator; under tensor hypercontraction (THC) the vertex factorises further, but the
denominator does not, and every previous attempt in this project handled the mismatch by
partitioning the particle--hole pairs and treating each region with a different expansion.
This note derives a route that removes the partition. The denominator is represented on a
grid of exponentials shared by every pole of the certified rational approximation to
$\eta^{(0)}$ and by every moment order, the powers of the denominator are recovered
exactly by a binomial expansion in the orbital energies, and a single measured scalar
residual bounds the relative spectral error of the resulting auxiliary-space objects at
\emph{every} order. Three candidate treatments of the higher powers are compared against an
exact evaluation; two are rejected on measurement. The cost model and the measured THC rank
requirement are given: every stage is cubic in operations and quadratic in memory, which
is the property the whole exercise exists to obtain, while the constants place the
crossover against the density-fitted route near one thousand orbitals.
\end{abstract}

\section*{The requirement}

Every stage of a tensor-hypercontracted path must be $O(N^3)$ in operations and $O(N^2)$
in memory. Nothing else justifies the construction: the density-fitted route is already
available, is simpler, and carries a certified error budget, so a hypercontracted route
that is merely faster by a constant has no reason to exist. What it buys is the exponent.

The distinction between constant and exponent is worth stating sharply, because the two
have different owners. The constant is dominated by the hypercontraction rank, a property
of the factorisation algorithms, which are an active field and are expected to improve.
The exponent is a property of the solver. No improvement in the integrals can repair a
quartic contraction. A poor constant is therefore a reason to wait; a poor exponent is a
defect, and the scaling of each stage is reported below on that basis rather than as an
afterthought.

\section{The object to be computed}

Write the particle--hole gaps as $D_{ia} = \epsilon_a - \epsilon_i > 0$ and let
$\mat{V} \in \mathbb{R}^{N_{\text{aux}} \times N_{ov}}$ be the vertex, so that the
occupied--virtual electron repulsion integrals are
\begin{align}
    (ia|jb) = \sum_P V_{P,ia} V_{P,jb}.
\end{align}
The moments of the density--density response, projected into the auxiliary space, are
\begin{align}
    \bar{\bm\eta}^{(m)} = \mat{V} \bm\eta^{(m)} \mat{V}^T,
    \qquad
    \bm\eta^{(m)} = (\mat{X}+\mat{Y})\, \bm\Omega^m (\mat{X}+\mat{Y})^T ,
\end{align}
with $\bm\Omega$ the RPA excitation energies and $\mat{X}$, $\mat{Y}$ the RPA
eigenvectors. These are the quantities the self-energy convolution consumes. The symbols
$\mat{X}$ and $\mat{Y}$ appear only in this equation; from Section~2 onwards $X$ denotes
the THC collocation matrix.

Two families of auxiliary-space objects generate all of them. The first is purely
polynomial in the gaps,
\begin{align}
    \Dcal^{(m)} = \mat{V}\, \mat{D}^m\, \mat{V}^T
    \qquad \text{i.e.} \qquad
    \Dcal^{(m)}_{PQ} = \sum_{ia} V_{P,ia}\, D_{ia}^m\, V_{Q,ia} ,
    \label{eq:Dcal}
\end{align}
and the second carries the RPA screening,
\begin{align}
    \Ncal^{(m)}
    = \frac{2}{\pi} \int_0^\infty
      \parth{\mat{I} + 4\Qcal^{(1)}(t)}^{-1} \Qcal^{(m+1)}(t) \; \mathrm{d}t ,
    \qquad
    \Qcal^{(k)}(t) = \mat{V}\, \frac{\mat{D}^{k}}{\mat{D}^2 + t^2}\, \mat{V}^T .
    \label{eq:compact}
\end{align}
Equation~\eqref{eq:compact} is the compact form of the RPA integral: it follows from
writing $\bm\eta^{(0)} = (\mat{A}+\mat{B})^{-1}[(\mat{A}+\mat{B})(\mat{A}-\mat{B})]^{1/2}$,
using $\mat{M}^{1/2} = \nicefrac{2}{\pi}\,\mat{M}\int_0^\infty (t^2\mat{I}+\mat{M})^{-1}
\mathrm{d}t$, applying the Woodbury identity, and then eliminating the bare term with the
push-through identity $(\mat{I}+\mat{M})^{-1} = \mat{I} - \mat{M}(\mat{I}+\mat{M})^{-1}$.
Both terms that would otherwise have to cancel against one another disappear analytically.

The two families combine into every moment through one recursion. Writing
\begin{align}
    \Dcal^{(m)} \;\text{for odd } m, \qquad \Ncal^{(m)} \;\text{for even } m,
\end{align}
and calling either $\mat{A}^{(m)}$,
\begin{align}
    \bar{\bm\eta}^{(m)}
    = \mat{A}^{(m)} + 4 \!\!\sum_{\substack{k = 1,3,5,\dots \\ k \le m-1}}\!\!
      \mat{A}^{(m-1-k)}\, \bar{\bm\eta}^{(k)} ,
    \qquad \bar{\bm\eta}^{(0)} = \Ncal^{(0)} ,\quad \bar{\bm\eta}^{(1)} = \Dcal^{(1)} .
    \label{eq:recursion}
\end{align}
The parity of $m-1-k$ always matches that of $m$, so the odd moments are built entirely
from $\Dcal$ and the even ones entirely from $\Ncal$. This is an $O(N_{\text{aux}}^3)$
recursion that never touches the particle--hole index.

\subsection{The rational form of the integral}

In production the $t$-integral in \eqref{eq:compact} is not quadratured. The inverse square
root is replaced by a certified rational approximation
$x^{-1/2} \approx \sum_j w_j (x + s_j)^{-1}$ on a rigorous spectral enclosure of
$\tilde{\mat{M}} = \mat{D}^2 + 4\mat{W}\mat{W}^T$, $\mat{W} = \mat{D}^{1/2}\mat{V}^T$, whose
poles $s_j$ and weights $w_j$ are chosen against a \emph{measured} scalar error. Equation
\eqref{eq:compact} becomes a finite sum,
\begin{align}
    \Ncal^{(m)}
    = \sum_j w_j \parth{\mat{I} + 4\Qcal^{(1)}(s_j)}^{-1} \Qcal^{(m+1)}(s_j) ,
    \qquad
    \Qcal^{(k)}(s) = \mat{V}\, \frac{\mat{D}^{k}}{\mat{D}^2 + s}\, \mat{V}^T ,
    \label{eq:poles}
\end{align}
with the $\nicefrac{2}{\pi}$ absorbed into $w_j$. At the production tolerance of $10^{-14}$
this is 18 to 19 poles on the systems used below. Everything that follows targets
\eqref{eq:poles}: the shifts are a fixed, known set rather than a grid whose scale has to be
optimised per calculation.

\section{Why tensor hypercontraction does not simply substitute}

Under THC the vertex factorises across the orbital pair,
\begin{align}
    (ia|jb) = \sum_{PQ} X_{iP} X_{aP}\, Z_{PQ}\, X_{jQ} X_{bQ} ,
\end{align}
so the natural replacement for $\mat{V}$ is $\mat{L}^T \mat{X}$, where
$X_{P,ia} = X_{iP} X_{aP}$ and $\mat{Z} = \mat{L}\mat{L}^T$. Substituting into
\eqref{eq:Dcal} costs nothing. Referencing the orbital energies to the chemical potential
$\mu$ --- writing $\eb_i = \epsilon_i - \mu$ and $\eb_a = \epsilon_a - \mu$, which leaves
$D_{ia} = \eb_a - \eb_i$ unchanged --- the power expands binomially into terms separable in
$i$ and $a$,
\begin{align}
    \Dcal^{(m)} \to \mat{L}^T \Dcal_X^{(m)} \mat{L} ,
    \qquad
    \Dcal_{X,PQ}^{(m)} = \sum_{p=0}^{m} \binom{m}{p}
        \sqrbrckts{\sum_i X_{iP} X_{iQ} (-\eb_i)^p}
        \sqrbrckts{\sum_a X_{aP} X_{aQ} \eb_a^{\,m-p}} ,
    \label{eq:binomialD}
\end{align}
which is $O(N^3)$ and is how the THC dTDA already works. The shift is not needed here ---
no exponentials are involved --- but it costs nothing and keeps the expansion identical to
the one Section~4 needs.

Substituting into \eqref{eq:poles} does not work, because $\mat{D}^k/(\mat{D}^2+s)$ is not
separable. Every earlier attempt in this project addressed that by splitting the
particle--hole pairs at each quadrature point $z$ into $D \ll z$, $D \sim z$ and $D \gg z$
regions, using the asymptotic series
$D^{2r+1}/(D^2+z^2) = \sum_c (-1)^c D^{2r+1+2c} z^{-2c-2}$ on one side, the Laplace/cosine
integral $\int_0^\infty \cos(zp)\, e^{-Dp}\,\mathrm{d}p$ on the other, and correction terms
across the boundary. The separable builders can only mask on $i$ and $a$ independently, so
the region --- a staircase in the sorted $(i,a)$ grid --- has to be approximated by unions of
rectangles, and the bookkeeping dominates.

That partition is \emph{not} required by the factorisation. It is required by the
oscillatory inner quadrature, which cannot resolve $\cos(zp)$ against $e^{-Dp}$ across a
wide range of $D$ at a single scale. Replacing the quadrature removes the partition.

\section{A shared exponential representation}

\subsection{The representation}

With the same referencing to $\mu$, so that $\eb_i < 0 < \eb_a$, on a set of positive
exponents $\{t_k\}$,
\begin{align}
    g_s(D) \equiv \frac{D}{D^2+s} \; \approx \; \sum_k c_k(s)\, e^{-t_k D} ,
    \qquad
    e^{-t_k D} = e^{t_k \eb_i}\, e^{-t_k \eb_a} .
    \label{eq:expfit}
\end{align}
The shift enters only through the coefficients: \textbf{the exponents are shared by every
pole and by every moment order}, so whatever a solver builds per exponent is built once.
The $\mu$ reference is load-bearing rather than cosmetic --- it puts both factors in
$[0,1]$, so neither overflows however wide the orbital range, and it makes the binomial in
\eqref{eq:binomialD} a sum of like-signed terms.

\subsection{What the residual certifies}

The coefficients are obtained by weighted least squares against the \emph{true} set of gaps,
minimising $\sum_{ia} \big[(\sum_k c_k e^{-t_k D_{ia}} - g_s(D_{ia}))/g_s(D_{ia})\big]^2$, so
the reported residual
\begin{align}
    \varepsilon = \max_{ia} \frac{|\delta(D_{ia})|}{g_s(D_{ia})} ,
    \qquad
    \delta(D) = \sum_k c_k e^{-t_k D} - g_s(D) ,
\end{align}
is realised on the system at hand rather than bounded over an enclosing interval. It then
transfers directly to the operator, and --- this is the point --- it does so at every moment
order. Since $\Delta \Qcal^{(m+1)} = \mat{V}\,\mathrm{diag}(\mat{D}^m \delta)\,\mat{V}^T$ is
symmetric, for any unit $\mat{u}$
\begin{align}
    \big| \mat{u}^T \Delta \Qcal^{(m+1)} \mat{u} \big|
    \;=\; \Big| \sum_{ia} D_{ia}^m\, \delta(D_{ia})\, (\mat{V}^T\mat{u})_{ia}^2 \Big|
    \;\le\; \varepsilon \sum_{ia} D_{ia}^m\, g_s(D_{ia})\, (\mat{V}^T\mat{u})_{ia}^2
    \;=\; \varepsilon\, \mat{u}^T \Qcal^{(m+1)} \mat{u} ,
\end{align}
because $D^m g_s(D) > 0$ makes $\Qcal^{(m+1)}$ positive semidefinite. Hence
\begin{align}
    \big\| \Delta \Qcal^{(m+1)} \big\|_2 \;\le\; \varepsilon\, \big\| \Qcal^{(m+1)} \big\|_2
    \qquad \text{for every } m .
    \label{eq:bound}
\end{align}
One scalar number certifies the whole moment set, and it does not degrade with order.

\subsection{Selecting the exponents}

The exponential basis on a real gap distribution has a numerical rank near 30 in double
precision: the least-squares effective rank saturates between 27 and 34 whether a geometric
grid holds 48 points or 120. A geometric layout therefore reaches the achievable floor only
by oversampling. Ranking a dense set of candidates by pivoted QR on the column-normalised
basis $[e^{-t D_{ia}}]$ and keeping the leading pivots reaches the same floor with far
fewer points (Table~\ref{tab:grid}).

\begin{table}[h]
\centering
\begin{tabular}{rcc}
\toprule
& \multicolumn{2}{c}{worst relative residual on $g_s$, all poles} \\
\cmidrule(l){2-3}
points & geometric & pivot-selected \\
\midrule
36 & $4.4\times10^{-8}$ & $1.4\times10^{-10}$ \\
42 & $4.5\times10^{-9}$ & $\bm{1.0\times10^{-12}}$ \\
60 & $1.9\times10^{-11}$ & $2.9\times10^{-12}$ \\
\bottomrule
\end{tabular}
\caption{Benzene, cc-pVDZ, 18 poles at a requested $\eta^{(0)}$ tolerance of $10^{-14}$.
Solver cost is linear in the number of exponents, so the saving is direct.}
\label{tab:grid}
\end{table}

Below a few times $10^{-13}$ the fit is limited by double precision rather than by the
grid, and more points do not help; a request below that is rejected rather than silently
answered. Table~\ref{tab:sizes} gives the sizes the search returns.

\begin{table}[h]
\centering
\begin{tabular}{lccc}
\toprule
& \multicolumn{3}{c}{exponents required} \\
\cmidrule(l){2-4}
system & $\varepsilon \le 10^{-6}$ & $10^{-9}$ & $10^{-11}$ \\
\midrule
lithium hydride & 16 & 19 & 21 \\
water           & 22 & 28 & 32 \\
ozone           & 23 & 28 & 36 \\
benzene         & 25 & 33 & 40 \\
\bottomrule
\end{tabular}
\caption{cc-pVDZ, 18 to 19 poles. The measured residual is the realised one on the
system's own gaps, and the operator error it certifies through \eqref{eq:bound} runs one to
two decades inside it.}
\label{tab:sizes}
\end{table}

\section{The higher powers}

Equation \eqref{eq:expfit} supplies $\Qcal^{(1)}$. The moments need $\Qcal^{(m+1)}$ for
even $m$ up to $n_{\max}-1$, and there are three ways to reach them.

\paragraph{A: polynomial division.} The identity
\begin{align}
    \frac{D^{2r+1}}{D^2+s}
    = \sum_{k=0}^{r-1} (-s)^k D^{2r-1-2k} \; + \; (-s)^r \frac{D}{D^2+s}
    \label{eq:division}
\end{align}
is exact for every pair, and every term but the last is a shift-independent $\Dcal^{(j)}$
that the dTDA already builds. Contracted into the auxiliary space this needs only two
builds per exponent and one set of $\Dcal$ for the whole calculation, so it is much the
cheapest. It is also, in disguise, what the previous implementations' \texttt{build\_large}
plus \texttt{build\_correction\_large} amounted to.

\paragraph{B: binomial in the orbital energies.} Keep the fit at $m=0$ and recover the
power exactly, as in \eqref{eq:binomialD}:
\begin{align}
    \Qcal_X^{(m+1)}(s) \;\approx\; \sum_k c_k(s) \sum_{p=0}^{m} \binom{m}{p}\;
      \mat{O}_p(t_k) \Had \mat{U}_{m-p}(t_k) ,
    \label{eq:routeB}
\end{align}
with $\Had$ the elementwise product and
\begin{align}
    \big[\mat{O}_p(t)\big]_{PQ} = \sum_i X_{iP} X_{iQ}\, (-\eb_i)^p\, e^{t \eb_i} ,
    \qquad
    \big[\mat{U}_q(t)\big]_{PQ} = \sum_a X_{aP} X_{aQ}\, \eb_a^{\,q}\, e^{-t \eb_a} .
    \label{eq:OU}
\end{align}
Every factor is non-negative in sign structure, so no cancellation is introduced. The cost
is $(n_{\max}+1)$ builds per exponent instead of two.

\paragraph{C: fit the target directly.} Fit $D^{m+1}/(D^2+s)$ on the same exponents and
keep two builds per exponent at every order.

\begin{table}[h]
\centering
\begin{tabular}{cccc}
\toprule
& \multicolumn{3}{c}{relative error on $\Ncal^{(m)}$, max norm} \\
\cmidrule(l){2-4}
$m$ & A: division & B: binomial & C: direct fit \\
\midrule
0 & $5.3\times10^{-13}$ & $5.3\times10^{-13}$ & $5.3\times10^{-13}$ \\
2 & $9.7\times10^{-11}$ & $9.7\times10^{-13}$ & $6.8\times10^{-8}$ \\
4 & $1.0\times10^{-8}$  & $1.0\times10^{-12}$ & $2.3\times10^{-4}$ \\
6 & $1.1\times10^{-6}$  & $1.0\times10^{-12}$ & $1.2\times10^{-1}$ \\
8 & $\bm{1.2\times10^{-4}}$ & $\bm{9.8\times10^{-13}}$ & $\bm{9.3\times10^{-1}}$ \\
\bottomrule
\end{tabular}
\caption{Benzene, cc-pVDZ, 42 exponents, measured against an exact evaluation of the same
pole sum \eqref{eq:poles} with a density-fitted vertex, so that only the treatment of the
denominator is under test.}
\label{tab:routes}
\end{table}

Route A loses about one and a half decades per moment order. The terms of
\eqref{eq:division} carry factors of $s^k$ that the answer does not, and although the
cancellation is benign where $s \lesssim D_{\max}^2$, it is not where the pole shift is much
larger. A finer grid does not recover it, because the loss is not in the fit.

Route C saturates: at 100 exponents it is no better than at 40. A function growing like
$D^{m-1}$ cannot be built from decaying exponentials without cancellation that double
precision will not carry.

Route B is flat in the moment order at $10^{-12}$, exactly as \eqref{eq:bound} predicts,
and is the route adopted. It is worth noting that this kernel --- the binomial expansion of
$D^m e^{-tD}$ --- is the one the earlier implementations already contained; what changes is
everything around it.

\section{Removing the Cholesky factor}

Substituting $\mat{V} \to \mat{L}^T\mat{X}$ into \eqref{eq:poles} would appear to require
$\mat{L}$, a Cholesky factor of $\mat{Z}$, since $\Qcal^{(k)} = \mat{L}^T \Qcal_X^{(k)}
\mat{L}$ with $\Qcal_X^{(k)} = \mat{X}[\mat{D}^k/(\mat{D}^2+s)]\mat{X}^T$. The push-through
identity $(\mat{I}+\mat{A}\mat{B})^{-1}\mat{A} = \mat{A}(\mat{I}+\mat{B}\mat{A})^{-1}$
removes it:
\begin{align}
    \parth{\mat{I} + 4\mat{L}^T \Qcal_X^{(1)} \mat{L}}^{-1} \mat{L}^T
    = \mat{L}^T \parth{\mat{I} + 4 \Qcal_X^{(1)} \mat{L}\mat{L}^T}^{-1}
    = \mat{L}^T \parth{\mat{I} + 4 \Qcal_X^{(1)} \mat{Z}}^{-1} ,
\end{align}
so that $\Ncal^{(m)} = \mat{L}^T \Ncal_X^{(m)} \mat{L}$ with
\begin{align}
    \Ncal_X^{(m)} = \sum_j w_j \parth{\mat{I} + 4\Qcal_X^{(1)}(s_j)\mat{Z}}^{-1}
                      \Qcal_X^{(m+1)}(s_j) .
    \label{eq:NX}
\end{align}
The price is a non-symmetric solve in place of a Cholesky one, at the same
$O(N_{\text{THC}}^3)$. The same substitution turns \eqref{eq:recursion} into
\begin{align}
    \bm\zeta^{(m)}
    = \mat{A}_X^{(m)} + 4 \!\!\sum_{\substack{k=1,3,5,\dots \\ k \le m-1}}\!\!
      \mat{A}_X^{(m-1-k)}\, \mat{Z}\, \bm\zeta^{(k)} ,
    \qquad \bar{\bm\eta}^{(m)} = \mat{L}^T \bm\zeta^{(m)} \mat{L} ,
    \label{eq:zeta}
\end{align}
with $\mat{A}_X^{(m)}$ equal to $\Dcal_X^{(m)}$ for odd $m$ and $\Ncal_X^{(m)}$ for even
$m$. The $\bm\zeta$ produced here is in the same convention the existing THC dTDA
self-energy builder consumes, which applies the two outer factors of $\mat{Z}$ at
contraction time; no interface change is needed downstream.

\section{Algorithm and cost}

\begin{algorithm}[h]
\caption{THC dRPA density--density moments}
\begin{algorithmic}[1]
\State Reference the orbital energies to $\mu$; form the gaps $D_{ia}$.
\State Bound the spectrum of $\tilde{\mat{M}}$; select poles $\{s_j, w_j\}$ against the
       requested $\eta^{(0)}$ tolerance.
\State Rank candidate exponents by pivoted QR; grow the grid until the measured residual
       $\varepsilon$ meets its tolerance; fit $c_k(s_j)$ for every pole. \hfill $O(N_{ov} n_t^2)$
\For{each exponent $t_k$}
    \State build $\mat{O}_p(t_k)$, $\mat{U}_q(t_k)$ for $p,q \le n_{\max}$ \hfill $O(n_{\max} N_{\text{THC}}^2 N)$
    \State $\mat{R}_m(t_k) \gets \sum_p \binom{m}{p} \mat{O}_p \Had \mat{U}_{m-p}$ for even $m$
\EndFor
\State $\Qcal_X^{(m+1)}(s_j) \gets \sum_k c_k(s_j) \mat{R}_m(t_k)$ \hfill $O(n_{\text{pole}} n_{\max} n_t N_{\text{THC}}^2)$
\State $\Ncal_X^{(m)} \gets \sum_j w_j (\mat{I} + 4\Qcal_X^{(1)}(s_j)\mat{Z})^{-1} \Qcal_X^{(m+1)}(s_j)$
\State $\Dcal_X^{(m)}$ for odd $m$ by \eqref{eq:binomialD}, without the exponential factors.
\State $\bm\zeta^{(m)}$ by the recursion \eqref{eq:zeta}. \hfill $O(n_{\max}^2 N_{\text{THC}}^3)$
\end{algorithmic}
\end{algorithm}

Counting the passes each side actually takes, the leading terms are
\begin{align}
    \text{separable:} \quad
        & n_t (n_{\max}+1) N_{\text{THC}}^2 N
          \;+\; \tfrac{1}{2}(n_{\max}+1)^2 N_{\text{THC}}^3 , \\
    \text{density-fitted:} \quad
        & (3 n_{\text{pole}} + 2 n_{\max})\, N_{\text{aux}}^2\, N_{ov} .
\end{align}
The separable route is cubic in the system size and the density-fitted route quartic, so
there is a crossover; where it falls is set almost entirely by the THC rank, which enters
squared. Table~\ref{tab:crossover} evaluates the model at benzene's shape ratios.

The self-energy stage that consumes $\bm\zeta^{(m)}$ has to meet the same requirement,
and in a direct implementation it does not. Written as a loop over the Green's function
orbitals it contracts an $(N_{\text{THC}}, N)$ vertex against $\bm\zeta$ once per
orbital, which is $O(N_{\text{THC}}^2 N^2)$ operations over an $O(N^3)$ array. The
orbital index reaches the output only through the weighted sums the moment convolution
takes, so moving that sum inside the auxiliary contraction gives
\begin{align}
    \Sigma^{(n)} = \sum_{t \le n} \binom{n}{t} (\pm 1)^t\;
        \mat{X}^T \parth{ \bm\zeta'^{(t)} \Had \mat{G}^{(n-t)} } \mat{X} ,
    \qquad
    G^{(j)}_{PQ} = \sum_{x \in \text{sector}} \epsilon_x^{\,j} X_{xP} X_{xQ} ,
\end{align}
with the sign negative for the occupied sector and the whole $t$ sum accumulated in the
auxiliary space before a single rotation per order. Neither the array nor the loop
survives, and the stage becomes $O(n_{\max}^2 (N_{\text{THC}}^2 N + N_{\text{THC}} N^2))$
operations in $O(n_{\max}(N_{\text{THC}}^2 + N^2))$ memory.

\begin{table}[h]
\centering
\begin{tabular}{lccc}
\toprule
& \multicolumn{3}{c}{crossover, orbitals} \\
\cmidrule(l){2-4}
THC rank & $n_t = 24$ & $n_t = 33$ & $n_t = 42$ \\
\midrule
$5 \times N$  & 114 & 134 & 155 \\
$8 \times N$  & 380 & 433 & 486 \\
$10 \times N$ & 686 & 769 & 853 \\
$12 \times N$ & 1121 & 1241 & 1361 \\
\bottomrule
\end{tabular}
\caption{$n_{\max} = 9$, 18 poles at an $\eta^{(0)}$ tolerance of $10^{-14}$, benzene's
auxiliary-to-orbital and pair-to-orbital ratios held fixed.}
\label{tab:crossover}
\end{table}

\section{Error budget}

Three approximations enter, and each is measured separately rather than inferred from the
others.

\begin{enumerate}
\item \textbf{The pole set.} Certified by the measured scalar error of the rational
approximation to $x^{-1/2}$ over a rigorous enclosure of $\tilde{\mat{M}}$; at the
production default this is $\sim\!10^{-15}$ on 18 to 19 poles.
\item \textbf{The separable representation.} Certified by $\varepsilon$ through
\eqref{eq:bound}, at every moment order. Measured, the realised operator error runs one to
two decades inside the bound.
\item \textbf{The factorisation itself.} Not certified --- measured. This is the dominant
term in practice and it is discussed below.
\end{enumerate}

The factorisation error should be quoted against the \emph{density-fitted} integrals, not
the exact ones: least-squares THC fits the density-fitted tensor, and on water in cc-pVDZ
the auxiliary basis contributes $4.4\times10^{-3}$ relative on its own, which is larger than
anything the factorisation adds and hides it completely.

Table~\ref{tab:rank} gives the measured relationship for the most efficient construction
available. The frontier shift is quoted for the worst of the systems tested at each accuracy
rather than the average, because the systems do not fail together: at a rank of
$8 \times N$ water and ozone are already inside 6~meV while lithium hydride is out at
88~meV, its diffuse virtual orbital being the first thing the fit loses.

\begin{table}[h]
\centering
\begin{tabular}{ccl}
\toprule
rank $/\,N$ & relative Frobenius error & worst frontier shift \\
\midrule
$4.6$ -- $5.7$   & $2.8\times10^{-2}$ -- $6.9\times10^{-2}$ & 278~meV \\
$7.3$ -- $8.9$   & $2.1\times10^{-3}$ -- $3.6\times10^{-3}$ & 88~meV \\
$8.1$ -- $10.7$  & $6.4\times10^{-4}$ -- $1.0\times10^{-3}$ & 7.7~meV \\
$8.7$ -- $11.4$  & $1.8\times10^{-4}$ -- $3.5\times10^{-4}$ & 2.6~meV \\
$9.0$ -- $13.6$  & $4.5\times10^{-5}$ -- $1.8\times10^{-4}$ & 1.1~meV \\
\bottomrule
\end{tabular}
\caption{Least-squares THC with randomly-pivoted Cholesky point selection, cc-pVDZ, across
water, lithium hydride and ozone; error against the density-fitted tensor in the same
auxiliary basis, frontier shift against the density-fitted dTDA at $n_{\max}=3$.}
\label{tab:rank}
\end{table}

A few meV therefore costs a rank near $11 \times N$, and the requirement does not fall with
system size: benzene at $N = 114$ reaches $-2.5$~meV at $9.96 \times N$, the same band as
the small systems. Read together with Table~\ref{tab:crossover}, that places the crossover
against the density-fitted route near a thousand orbitals. The exponent grid is already
close to its own floor, so neither of the two levers has much left to give.

\section{Summary}

The particle--hole denominator separates cleanly once it is represented on a shared
exponential grid fitted per pole against the true gaps, with the powers of the denominator
recovered exactly by a binomial expansion in the orbital energies rather than fitted. The
region partitioning that dominated earlier implementations is an artefact of the oscillatory
inner quadrature, not of the factorisation, and it disappears with it. A single measured
scalar residual certifies the relative spectral error of every auxiliary-space object at
every moment order, and the Cholesky factor of the THC Coulomb matrix is removable by a
push-through identity.

Every stage is then cubic in operations and quadratic in memory, which is the property the
construction exists to obtain.

The constants are another matter. Solver cost scales with the square of the
hypercontraction rank; the rank needed for few-meV quasiparticle agreement is near
$11 \times N$ and does not fall with system size; and the resulting crossover against the
density-fitted route sits near one thousand orbitals. That is a statement about today's
factorisation algorithms rather than about this construction, and it is the half of the
problem that improvement in the integrals can reach. The exponent is the half that it
cannot, and it is now correct.

\end{document}
